% Options for packages loaded elsewhere
\PassOptionsToPackage{unicode}{hyperref}
\PassOptionsToPackage{hyphens}{url}
\documentclass[
]{article}
\usepackage{xcolor}
\usepackage[margin=1in]{geometry}
\usepackage{amsmath,amssymb}
\setcounter{secnumdepth}{-\maxdimen} % remove section numbering
\usepackage{iftex}
\ifPDFTeX
  \usepackage[T1]{fontenc}
  \usepackage[utf8]{inputenc}
  \usepackage{textcomp} % provide euro and other symbols
\else % if luatex or xetex
  \usepackage{unicode-math} % this also loads fontspec
  \defaultfontfeatures{Scale=MatchLowercase}
  \defaultfontfeatures[\rmfamily]{Ligatures=TeX,Scale=1}
\fi
\usepackage{lmodern}
\ifPDFTeX\else
  % xetex/luatex font selection
\fi
% Use upquote if available, for straight quotes in verbatim environments
\IfFileExists{upquote.sty}{\usepackage{upquote}}{}
\IfFileExists{microtype.sty}{% use microtype if available
  \usepackage[]{microtype}
  \UseMicrotypeSet[protrusion]{basicmath} % disable protrusion for tt fonts
}{}
\makeatletter
\@ifundefined{KOMAClassName}{% if non-KOMA class
  \IfFileExists{parskip.sty}{%
    \usepackage{parskip}
  }{% else
    \setlength{\parindent}{0pt}
    \setlength{\parskip}{6pt plus 2pt minus 1pt}}
}{% if KOMA class
  \KOMAoptions{parskip=half}}
\makeatother
\usepackage{color}
\usepackage{fancyvrb}
\newcommand{\VerbBar}{|}
\newcommand{\VERB}{\Verb[commandchars=\\\{\}]}
\DefineVerbatimEnvironment{Highlighting}{Verbatim}{commandchars=\\\{\}}
% Add ',fontsize=\small' for more characters per line
\usepackage{framed}
\definecolor{shadecolor}{RGB}{248,248,248}
\newenvironment{Shaded}{\begin{snugshade}}{\end{snugshade}}
\newcommand{\AlertTok}[1]{\textcolor[rgb]{0.94,0.16,0.16}{#1}}
\newcommand{\AnnotationTok}[1]{\textcolor[rgb]{0.56,0.35,0.01}{\textbf{\textit{#1}}}}
\newcommand{\AttributeTok}[1]{\textcolor[rgb]{0.13,0.29,0.53}{#1}}
\newcommand{\BaseNTok}[1]{\textcolor[rgb]{0.00,0.00,0.81}{#1}}
\newcommand{\BuiltInTok}[1]{#1}
\newcommand{\CharTok}[1]{\textcolor[rgb]{0.31,0.60,0.02}{#1}}
\newcommand{\CommentTok}[1]{\textcolor[rgb]{0.56,0.35,0.01}{\textit{#1}}}
\newcommand{\CommentVarTok}[1]{\textcolor[rgb]{0.56,0.35,0.01}{\textbf{\textit{#1}}}}
\newcommand{\ConstantTok}[1]{\textcolor[rgb]{0.56,0.35,0.01}{#1}}
\newcommand{\ControlFlowTok}[1]{\textcolor[rgb]{0.13,0.29,0.53}{\textbf{#1}}}
\newcommand{\DataTypeTok}[1]{\textcolor[rgb]{0.13,0.29,0.53}{#1}}
\newcommand{\DecValTok}[1]{\textcolor[rgb]{0.00,0.00,0.81}{#1}}
\newcommand{\DocumentationTok}[1]{\textcolor[rgb]{0.56,0.35,0.01}{\textbf{\textit{#1}}}}
\newcommand{\ErrorTok}[1]{\textcolor[rgb]{0.64,0.00,0.00}{\textbf{#1}}}
\newcommand{\ExtensionTok}[1]{#1}
\newcommand{\FloatTok}[1]{\textcolor[rgb]{0.00,0.00,0.81}{#1}}
\newcommand{\FunctionTok}[1]{\textcolor[rgb]{0.13,0.29,0.53}{\textbf{#1}}}
\newcommand{\ImportTok}[1]{#1}
\newcommand{\InformationTok}[1]{\textcolor[rgb]{0.56,0.35,0.01}{\textbf{\textit{#1}}}}
\newcommand{\KeywordTok}[1]{\textcolor[rgb]{0.13,0.29,0.53}{\textbf{#1}}}
\newcommand{\NormalTok}[1]{#1}
\newcommand{\OperatorTok}[1]{\textcolor[rgb]{0.81,0.36,0.00}{\textbf{#1}}}
\newcommand{\OtherTok}[1]{\textcolor[rgb]{0.56,0.35,0.01}{#1}}
\newcommand{\PreprocessorTok}[1]{\textcolor[rgb]{0.56,0.35,0.01}{\textit{#1}}}
\newcommand{\RegionMarkerTok}[1]{#1}
\newcommand{\SpecialCharTok}[1]{\textcolor[rgb]{0.81,0.36,0.00}{\textbf{#1}}}
\newcommand{\SpecialStringTok}[1]{\textcolor[rgb]{0.31,0.60,0.02}{#1}}
\newcommand{\StringTok}[1]{\textcolor[rgb]{0.31,0.60,0.02}{#1}}
\newcommand{\VariableTok}[1]{\textcolor[rgb]{0.00,0.00,0.00}{#1}}
\newcommand{\VerbatimStringTok}[1]{\textcolor[rgb]{0.31,0.60,0.02}{#1}}
\newcommand{\WarningTok}[1]{\textcolor[rgb]{0.56,0.35,0.01}{\textbf{\textit{#1}}}}
\usepackage{graphicx}
\makeatletter
\newsavebox\pandoc@box
\newcommand*\pandocbounded[1]{% scales image to fit in text height/width
  \sbox\pandoc@box{#1}%
  \Gscale@div\@tempa{\textheight}{\dimexpr\ht\pandoc@box+\dp\pandoc@box\relax}%
  \Gscale@div\@tempb{\linewidth}{\wd\pandoc@box}%
  \ifdim\@tempb\p@<\@tempa\p@\let\@tempa\@tempb\fi% select the smaller of both
  \ifdim\@tempa\p@<\p@\scalebox{\@tempa}{\usebox\pandoc@box}%
  \else\usebox{\pandoc@box}%
  \fi%
}
% Set default figure placement to htbp
\def\fps@figure{htbp}
\makeatother
\setlength{\emergencystretch}{3em} % prevent overfull lines
\providecommand{\tightlist}{%
  \setlength{\itemsep}{0pt}\setlength{\parskip}{0pt}}
\usepackage{bookmark}
\IfFileExists{xurl.sty}{\usepackage{xurl}}{} % add URL line breaks if available
\urlstyle{same}
\hypersetup{
  pdftitle={Weakly Informative Priors},
  hidelinks,
  pdfcreator={LaTeX via pandoc}}

\title{Weakly Informative Priors}
\author{}
\date{\vspace{-2.5em}2026-04-17}

\begin{document}
\maketitle

\subsection{Introduction}\label{introduction}

Weakly informative priors are the modern default in applied Bayesian
work. They assign non-trivial probability to a plausible range of
parameter values while remaining broad enough that the data dominate the
posterior. They differ from ``non-informative'' priors (which can be
improper or pathological) and from informative priors (which encode
substantive knowledge): the goal is to regularise the inference, rule
out absurd values, and stabilise sampling without injecting domain
claims that the data cannot challenge. The result is a posterior that
looks much like a frequentist estimate when data are abundant, but that
behaves sensibly under separation, sparsity, and high-dimensional
regression where likelihood-only methods misbehave.

\subsection{Prerequisites}\label{prerequisites}

A working understanding of priors, the role of the prior in Bayesian
updating, and basic familiarity with regression on standardised
predictors.

\subsection{Theory}\label{theory}

A prior is \emph{weakly informative} if it gives small but non-zero
probability to extreme values, has a finite second moment when possible,
and is centred at a value that does not bias the analysis (zero, in most
regression contexts). Standard choices for standardised predictors:

\begin{itemize}
\tightlist
\item
  Regression slopes: \(\mathcal N(0, 2.5)\) or \(\mathcal N(0, 5)\) on
  standardised covariates.
\item
  Intercept (after centring \(y\)): \(\mathcal N(0, 10)\).
\item
  Standard deviations of random effects or residuals:
  \(\mathrm{half\text{-}Normal}(0, 2.5)\),
  \(\mathrm{half\text{-}Cauchy}(0, 5)\), or
  \(\mathrm{Exponential}(\lambda)\) with rate matched to the SD scale.
\item
  Correlation matrices: LKJ prior with shape \(\eta \geq 2\), which
  gently regularises toward weakly correlated structures.
\end{itemize}

A useful calibration test: simulate from the prior predictive
distribution and check that the implied range of \texttt{y} is
plausible. If the prior implies effects of \(\pm 1{,}000\) on a quantity
bounded by physiology, it is too wide.

\subsection{Assumptions}\label{assumptions}

Predictors are on a sensible scale (typically standardised), priors are
proper (so the posterior is integrable), and the prior predictive range
covers the plausible space without dominating it.

\subsection{R Implementation}\label{r-implementation}

\begin{Shaded}
\begin{Highlighting}[]
\FunctionTok{library}\NormalTok{(brms)}

\CommentTok{\# Default brms priors are weakly informative}
\NormalTok{data\_ex }\OtherTok{\textless{}{-}} \FunctionTok{data.frame}\NormalTok{(}\AttributeTok{y =} \FunctionTok{rnorm}\NormalTok{(}\DecValTok{100}\NormalTok{), }\AttributeTok{x =} \FunctionTok{rnorm}\NormalTok{(}\DecValTok{100}\NormalTok{))}
\FunctionTok{get\_prior}\NormalTok{(y }\SpecialCharTok{\textasciitilde{}}\NormalTok{ x, }\AttributeTok{data =}\NormalTok{ data\_ex)}

\CommentTok{\# Explicit weakly informative}
\NormalTok{fit }\OtherTok{\textless{}{-}} \FunctionTok{brm}\NormalTok{(y }\SpecialCharTok{\textasciitilde{}}\NormalTok{ x, }\AttributeTok{data =}\NormalTok{ data\_ex,}
           \AttributeTok{prior =} \FunctionTok{prior}\NormalTok{(}\FunctionTok{normal}\NormalTok{(}\DecValTok{0}\NormalTok{, }\FloatTok{2.5}\NormalTok{), }\AttributeTok{class =} \StringTok{"b"}\NormalTok{) }\SpecialCharTok{+}
                   \FunctionTok{prior}\NormalTok{(}\FunctionTok{normal}\NormalTok{(}\DecValTok{0}\NormalTok{, }\DecValTok{5}\NormalTok{),   }\AttributeTok{class =} \StringTok{"Intercept"}\NormalTok{) }\SpecialCharTok{+}
                   \FunctionTok{prior}\NormalTok{(}\FunctionTok{exponential}\NormalTok{(}\DecValTok{1}\NormalTok{), }\AttributeTok{class =} \StringTok{"sigma"}\NormalTok{),}
           \AttributeTok{iter =} \DecValTok{2000}\NormalTok{, }\AttributeTok{chains =} \DecValTok{4}\NormalTok{, }\AttributeTok{refresh =} \DecValTok{0}\NormalTok{)}
\end{Highlighting}
\end{Shaded}

\subsection{Output \& Results}\label{output-results}

\texttt{get\_prior()} lists every parameter class together with the
package's default prior; explicit overrides replace them by class. After
fitting, \texttt{prior\_summary(fit)} confirms exactly which priors
entered the model --- a step worth running for any analysis intended for
publication.

\subsection{Interpretation}\label{interpretation}

A reporting sentence: ``Weakly informative priors \(\mathcal N(0, 2.5)\)
on standardised slopes and \(\mathrm{Exponential}(1)\) on the residual
scale yielded a posterior centred close to the OLS estimate
(\(\beta = 0.51\), 95 \% credible interval 0.31 to 0.71), with
negligible prior pull at this sample size; a sensitivity analysis under
\(\mathcal N(0, 1)\) produced essentially identical posterior
summaries.'' Disclosing both the prior and a sensitivity check is the
gold-standard report.

\subsection{Practical Tips}\label{practical-tips}

\begin{itemize}
\tightlist
\item
  Gelman's heuristic: half-Cauchy with scale 2.5--5 for variance
  components --- heavy enough to allow large variances when warranted,
  regularising enough to prevent boundary-zero estimates.
\item
  For unstandardised predictors, multiply the prior scale by the
  predictor's empirical SD to keep the implied effect range plausible.
\item
  LKJ(2) for correlation matrices is a sensible default; tighten only
  when you have reason to expect weak correlations.
\item
  Always run prior predictive checks
  (\texttt{sample\_prior\ =\ "only"}); they reveal absurd implied ranges
  before the data can mask them.
\item
  For non-standardised log-odds problems where the baseline log-odds is
  far from zero, widen the prior on the intercept
  (e.g.~\(\mathcal N(0, 10)\)) so the prior mode does not exclude the
  data-implied baseline.
\end{itemize}

\subsection{R Packages Used}\label{r-packages-used}

\texttt{brms} for formula-based prior specification with
\texttt{get\_prior()} / \texttt{prior\_summary()}, \texttt{rstanarm} for
an alternative interface with auto-scaled defaults, \texttt{bayestestR}
for posterior summaries that report the prior alongside the posterior,
and \texttt{bayesplot} for prior predictive visualisation.

\subsection{For Reviewers}\label{for-reviewers}

\textbf{What to look for in a paper using this method.}

\begin{itemize}
\tightlist
\item
  \textbf{Common misapplications.}
\item
  \textbf{Diagnostics that should be reported but often aren't.}
\item
  \textbf{Red flags in tables and figures.}
\item
  \textbf{What to verify.}
\item
  \textbf{What an adequate Methods paragraph must contain.}
\end{itemize}

\subsection{See also --- labs in this
chapter}\label{see-also-labs-in-this-chapter}

\begin{itemize}
\tightlist
\item
  \href{labs/lab-bayesian-thinking-control-vs-decision-error.Rmd}{Bayesian
  thinking; α control vs decision error}
\item
  \href{labs/lab-brms-stan-loo-hierarchical-models.Rmd}{brms / Stan,
  LOO, hierarchical models}
\end{itemize}

\end{document}
